% section: abstract
\begin{abstract}
Chart question answering is close to saturated: a current 2-billion-parameter vision-language model
scores \result{79.1} on ChartQA with no task-specific training, and the best published improvement
over a comparable prompted model is $+3.8$ points obtained on four H100 GPUs. Chart \emph{grounding}
--- identifying which region of the image an answer came from --- is not: the strongest published
result on RefChartQA's human subset is \result{32.83} AP@0.5, from a larger model trained one epoch
with boxes emitted as plain text.

We fine-tune one small vision-language model to emit, for each chart and question, a single strict
JSON record containing evidence regions, a typed expression tree, and a direct answer; a
deterministic interpreter then recomputes the answer from the expression tree. Because the
arithmetic leaves the model, reading the chart and computing with it become separately observable,
which supports an error decomposition ordinary systems cannot produce.

\TODO{results sentence: both headline before/after numbers with intervals, and the structured-output
cost, stated plainly including if negative}

We also report a mechanical account of why chart grounding is hard. The model represents the image
as a grid of $32\times32$-pixel visual tokens, and at the input resolution used, \result{53.2\%} of
RefChartQA grounding targets are narrower than one token on at least one axis --- they cannot be
resolved as distinct regions in principle. Grounding results are therefore stratified by target
size, where the aggregate figure conceals two different problems.

All training uses free-tier hardware and no paid compute. Every methodological choice is recorded in
an append-only decision log, and all decisions were frozen in a committed pre-registration before
any test split was opened.
\end{abstract}
